\chapter{Partial Circulants: A Running Laboratory}
\label{organized:ptc}
\label{chapter-ptc-laboratory}
\paragraph{Scope: an application of the general theory.}
The partial-circulant family is an application of the general transformation
semigroup and operator theory, not an exception to it. The additional
assumptions used for road-coloring questions are not imposed in this
chapter. The number of generators and their probabilities are specified
locally. In the associated stochastic chain, loops, coincident outgoing
targets, transient states or periodicity can occur. Stationary and limiting
statements use the hypotheses stated with each result; road-coloring
hypotheses must not be imported implicitly.

\section{Definitions and statement}
Let $\mathcal N=\{1,\ldots,n\}$. For $1\le i\le n$, define
\[
 xM_i=\begin{cases}i,&x=1,\\x-1,&x>1.\end{cases}
\]
We use right actions: $x(fg)=(xf)g$. Equivalently, transformation matrices
have a $1$ in position $(x,xf)$, and composition agrees with ordinary
matrix multiplication. Their matrix rank is the cardinality of their image.

Fix a nonempty set $I\subseteq\{1,\ldots,n\}$ and write
\[
 S=\langle M_i:i\in I\rangle,\qquad
 m=\max I,\qquad d=\gcd\{i:i\in I\},\qquad C=\{1,\ldots,m\}.
\]
Here $S$ consists of nonempty words in the generators. The \emph{kernel}
of a finite semigroup means its unique minimal nonempty two-sided ideal;
it is not to be confused with the fiber partition of a transformation.

For $r\in\mathbb{Z}/d\mathbb{Z}$, put
\[
 P_r=\{x\in\mathcal N:x\equiv r\pmod d\},\qquad C_r=P_r\cap C,
\]
and let
\[
 \mathcal A=\{A\subseteq C:|A\cap C_r|=1\text{ for every }r\in\mathbb{Z}/d\mathbb{Z}\}.
\]
For $A\in\mathcal A$, denote its unique point in $C_r$ by $a_r$. Define
\[
 xk_{A,t}=a_{[x]-t}\qquad(x\in\mathcal N,\ t\in\mathbb{Z}/d\mathbb{Z}),
\]
where $[x]$ denotes the residue class of $x$ modulo $d$.

\begin{theorem}\label{ptc-kernel:thm:main}
The following statements hold.
\begin{enumerate}
\item The minimum rank of an element of $S$ is $d$.
\item The kernel of $S$ is exactly
\[
 K=\{s\in S:\rank(s)=d\}
   =\{k_{A,t}:A\in\mathcal A,\ t\in\mathbb{Z}/d\mathbb{Z}\}.
\]
\item Every element of $K$ has fiber partition $\{P_r:r\in\mathbb{Z}/d\mathbb{Z}\}$.
Its possible ranges are exactly the members of $\mathcal A$.
\item Multiplication is given by
\[
 k_{A,t}k_{B,u}=k_{B,t+u}.
\]
Consequently $K\cong C_d\times R_{\mathcal A}$, where $C_d$ is the cyclic
group of order $d$ and $R_{\mathcal A}$ is the right-zero semigroup
on $\mathcal A$, with product $AB=B$. In particular, $K$ is a right group.
\item The number of ranges and the cardinality of the kernel are
\[
 |\mathcal A|=(m/d)^d,\qquad |K|=d(m/d)^d.
\]
\end{enumerate}
\end{theorem}


\section{Congruence classes and elementary collapses}

\begin{lemma}\label{ptc-kernel:lem:lower}
If $w$ is a word of length $\ell$, then
\[
 xw\equiv x-\ell\pmod d\qquad(x\in\mathcal N).
\]
Moreover $Cw\subseteq C$ and $|Cw|\ge d$. In particular, $\rank(w)\ge d$.
\end{lemma}
\begin{proof}
For every $i\in I$, we have $xM_i\equiv x-1\pmod d$: at $x=1$ this
uses $d\mid i$, and at all other points it follows from the definition.
Induction proves the formula for words. Since every generator preserves
$C$, so does every word. The set $C$ meets every residue class modulo
$d$, and the congruence formula shows that its image does as well.
\end{proof}

\begin{lemma}\label{ptc-kernel:lem:collapse}
For each $i\in I$ with $i<m$, and each $a\in C$, the restriction
semigroup $S|_C$ contains the transformation fixing every point other
than $a$ and sending $a$ to $a+i$ modulo $m$.
\end{lemma}
\begin{proof}
On $C$, let $c=M_m|_C$. This is the permutation $x\mapsto x-1$
modulo $m$, so $c^{-1}=c^{m-1}$. The restriction of
$M_iM_m^{m-1}$ fixes every $x>1$ and sends $1$ to $i+1$.
Conjugation by powers of $c$ gives all the asserted transformations.
These conjugations are realized by words in the original generators,
since both $c$ and $c^{-1}$ are restrictions of positive powers of $M_m$.
\end{proof}

\begin{lemma}\label{ptc-kernel:lem:retraction}
For every $A\in\mathcal A$, there exists $v_A\in S$ such that
\[
 xv_A=a_r\quad(x\in C_r),\qquad
 xv_A\equiv x\pmod d\quad(x\in\mathcal N).
\]
\end{lemma}
\begin{proof}
Form the directed graph on $\mathbb{Z}/m\mathbb{Z}$ with edges
\[
 a\longrightarrow a+i\qquad(i\in I,\ i<m).
\]
Its strongly connected components are exactly the classes modulo $d$.
Indeed, the increments generate the subgroup $d\mathbb{Z}/m\mathbb{Z}$. In a finite
cyclic group the additive inverse of any increment is a nonnegative
multiple of that increment, so directed reachability gives precisely
the cosets of this subgroup.

In each component choose a directed spanning tree oriented toward the
prescribed root $a_r$. Such a tree exists, for example by choosing from
each nonroot vertex an edge along a shortest directed path to the root.
For each nonroot vertex apply the elementary collapse to its parent,
in decreasing order of distance from the root. A point first moves to
its parent and is subsequently moved along the remaining path to the
root. Thus the product of these collapses is the desired retraction
on $C$.

Choose words realizing these elementary collapses as in
Lemma~\ref{ptc-kernel:lem:collapse}, and concatenate them to obtain $v_A$.
Each collapse preserves residues modulo $d$. By
Lemma~\ref{ptc-kernel:lem:lower}, a representing word for such a collapse must
have length divisible by $d$ (apply the congruence formula to any point
of $C$). Hence that word preserves residues on all of $\mathcal N$,
not only on $C$. The product $v_A$ does too.

If there are no nonroot vertices, then $d=m$, $I=\{m\}$, and $A=C$.
In this case take $v_A=M_m^m$, which has the required properties.
\end{proof}


\section{Proof of the theorem}
\begin{proof}[Proof of Theorem~\ref{ptc-kernel:thm:main}]
Choose a positive integer $L$ divisible by $d$ and satisfying $L\ge n-m$.
Iteration of $M_m$ moves every point outside $C$ down by one until it
enters $C$, and subsequently stays inside $C$. Therefore
\[
 \mathcal N M_m^L\subseteq C.
\]
For $A\in\mathcal A$ and an integer representative $0\le t<d$,
Lemma~\ref{ptc-kernel:lem:retraction} now gives
\[
 M_m^{L+t}v_A=k_{A,t}.
\]
Indeed the prefix sends $x$ to a point of $C$ congruent to $x-t$
modulo $d$, and $v_A$ sends that point to $a_{[x]-t}$.
Thus every $k_{A,t}$ belongs to $S$ and has rank $d$.
Lemma~\ref{ptc-kernel:lem:lower} proves that this is the minimum rank.

Conversely, suppose $w\in S$ has rank $d$, and choose a word
representing it of length $\ell$. Since $Cw\subseteq\mathcal Nw$ and
$|Cw|\ge d$, we have
\[
 \mathcal Nw=Cw\subseteq C.
\]
Its image therefore consists of exactly one point in each residue
class modulo $d$. Call this transversal $A$. The congruence formula
forces
\[
 xw=a_{[x]-[\ell]},
\]
so $w=k_{A,[\ell]}$. This proves the classification of all minimum-rank
elements and their fiber partitions and ranges.

For multiplication, observe that $a_r\equiv r\pmod d$, whence
\[
 xk_{A,t}k_{B,u}=a_{[x]-t}k_{B,u}=b_{[x]-t-u}.
\]
Thus $k_{A,t}k_{B,u}=k_{B,t+u}$. The parametrization is injective:
the image determines $A$, and the induced shift of residue classes
determines $t$. It follows that $(t,A)\mapsto k_{A,t}$ is an
isomorphism $C_d\times R_{\mathcal A}\to K$.

It remains to check that $K$ is the minimal ideal, rather than merely
a subsemigroup. Multiplication cannot increase transformation rank,
and no element of $S$ has rank below $d$. Hence $K$ is a two-sided
ideal. The multiplication formula also gives $kK=K$ for every $k\in K$.
If $J$ is any nonempty two-sided ideal, choose $s\in J$ and $k\in K$.
Then $sk\in J\cap K$, and
\[
 K=(sk)K\subseteq J.
\]
Thus $K$ is the unique minimal nonempty two-sided ideal.

Finally, each of the $d$ sets $C_r$ has $m/d$ points. Choosing one
point independently from each gives $|\mathcal A|=(m/d)^d$, and
$|K|=d|\mathcal A|$.
\end{proof}


\section{Further structure}
\begin{corollary}
The idempotents of $K$ are precisely $k_{A,0}$ for $A\in\mathcal A$.
They are the retractions onto the possible ranges. For each fixed $A$,
\[
 H_A=\{k_{A,t}:t\in\mathbb{Z}/d\mathbb{Z}\}
\]
is a maximal subgroup, cyclic of order $d$, with identity $k_{A,0}$.
There is one Green $\mathcal R$-class in $K$, while its Green
$\mathcal L$-classes (and hence its $\mathcal H$-classes) are the sets
$H_A$ indexed by the ranges.
\end{corollary}
\begin{proof}
All assertions follow from the direct-product multiplication.
In particular, $k_{A,t}^2=k_{A,2t}$ equals $k_{A,t}$ precisely when
$t=0$, and fixing $A$ gives the group law of $\mathbb{Z}/d\mathbb{Z}$.
\end{proof}

\begin{corollary}
The semigroup $S$ contains a constant transformation if and only if
$d=1$. In that case its kernel is the right-zero semigroup consisting
of the $m$ constant maps with values in $C$.
\end{corollary}
\begin{proof}
A constant transformation has rank one. Apply the minimum-rank formula
and the classification of kernel ranges.
\end{proof}

\begin{remark}[Composition convention]
All right-group statements above use $x(fg)=(xf)g$. If transformations
are multiplied with the opposite convention, the resulting semigroup
is the opposite semigroup: the corresponding kernel is a left group,
with the same rank, fiber partitions, and ranges.
\end{remark}



\section{The ladder chain and its stationary distribution}
Assign positive probabilities $p_i$ to the selected generators, with
$\sum_{i\in I}p_i=1$, and put $T=\sum_{i\in I}p_iM_i$.
Write $\mu=\sum_{i\in I}i p_i$. Then
\begin{proposition}\label{ptc-book:stationary}
The stationary distribution, extended by zero outside $C$, is
\[
 \pi_j=\frac{\sum_{i\in I,\ i\ge j}p_i}{\mu}\quad(1\le j\le m),
 \qquad \pi_j=0\quad(j>m).
\]
On $C$ the chain is irreducible and its period is $d=\gcd(I)$.
\end{proposition}
\begin{proof}
The stationary equations read $\pi_j=\pi_1p_j+\pi_{j+1}$ on the core,
where $p_j=0$ for unselected indices and $\pi_{m+1}=0$. Backward substitution
and normalization give the formula. The $M_m$ cycle visits all of $C$, so
the chain is irreducible there. A return from state 1 to itself consists of
excursions of lengths in $I$, and each selected length occurs. Their greatest
common divisor is therefore the period. States above $m$ descend into $C$.
\end{proof}


\section{The spectrum of the collision operator}
For a kernel-valued random transformation $K$, use the convention
\[
 N=\mathbb E[KK^{\mathsf T}],\qquad N_0=J-N.
\]
Every kernel transformation has the same partition, so $KK^{\mathsf T}$
is independent of $K$. Its $(x,y)$ entry is one exactly when
$x\equiv y\pmod d$. Consequently, after ordering the core by residue classes,
\[
 N|_{C\times C}=I_d\otimes J_q,\qquad q=m/d.
\]
The nonzero eigenvalue on the core is $q$, of multiplicity $d$; zero has
multiplicity $m-d$. The matrix $N/q$ restricted to the core is the orthogonal
projection onto vectors constant on each residue class. On the full state
space the nonzero eigenvalues are the full residue-block sizes, which need
not equal $q$ when transient states are present.


\section{Why permanental and common-letter hierarchies agree here}
For a matrix $X$, define $X^{\vee k}$ by permanents of $k$-by-$k$
submatrices, indexed by subsets. In general the permanental lift of an
average is not the average of the lifts. Here, however,
\begin{proposition}\label{ptc-book:permanents}
At every subset level,
\[
 T^{\vee k}=\sum_{i\in I}p_i M_i^{\vee k}.
\]
\end{proposition}
\begin{proof}
All generators differ only in row 1. A submatrix either omits that row,
in which case it is identical for every generator, or includes it once.
A permanent is linear in each row separately. In either case averaging
commutes with the submatrix permanent.
\end{proof}
Thus the common-letter semigroup hierarchy is also the permanental
hierarchy of the single ladder transition matrix. The equality is special
to this shared-row structure, not a general identity for stochastic matrices.
Identifying compatible fields, projection normalizations, and cyclic-pattern
recurrent classes is the next extension of this chapter, using the separate
hierarchy paper and the operator draft.


\section{The rank condition: a hypothesis test}
For distinct $a,b$, take $R=M_a$, $B=M_b$, on their irreducible core. Then
\[
 \Delta=\tfrac12e_1(e_a^{\mathsf T}-e_b^{\mathsf T}),
 \quad \pi\Delta=\tfrac{\pi_1}{2}(e_a^{\mathsf T}-e_b^{\mathsf T}).
\]
Since $\pi_1>0$, the row span generated by $\pi\Delta$ already equals
the one-dimensional row space of $\Delta$. This is Chapter 5's rank
condition. When $\gcd(a,b)=1$, the core is aperiodic and the synchronization
criterion applies. When the gcd exceeds one, the same rank condition holds
but synchronization is impossible. Thus aperiodicity is indispensable in
the surrounding criterion; a theorem statement separated from its standing
hypotheses would be misleading.
The example $a=2,b=4$ should be used as a routine hypothesis check.


\section{A worked example and exercises}
Take $n=m=6$, $I=\{2,4,6\}$, and equal probabilities. Then
\[
 d=2,\quad q=3,\quad
 \pi=(1/4,1/4,1/6,1/6,1/12,1/12).
\]
The partition is $(135\mid246)$, there are nine ranges and eighteen
kernel elements, and $N$ has eigenvalues $3,3,0,0,0,0$.
This chain is periodic, so ordinary powers need not converge; use
Cesaro or Abel projections.
\begin{enumerate}
\item Verify each stationary equation and that each partition block has
stationary mass $1/2$.
\item List the nine ranges and verify the kernel multiplication law.
\item Diagonalize the two $J_3$ blocks and interpret their eigenspaces.
\item Extend the example to $n=8$ without changing $I$. Determine the
transient states and the full-space spectrum of $N$.
\item Compare the common-letter and permanental level-two matrices in
this example. Find a non-PtC pair for which the equality fails.
\item Work out the rank-one case $I=\{2,5\}$ on five states.
\end{enumerate}


\section{Research directions, not asserted theorems}
The Lie-constant converse in the separate Lie notes remains open in this
project. Exact computations suggest that the Lie-constant intersection on
the recurrent core has dimension $m-d$; no proof for arbitrary PtC supports
is supplied here. Unicycle enumeration, canonical recoloring, and the double
construction are further tests of the general theory, with their original
hypotheses and normalization conventions kept explicit.

